\documentclass{article}
\usepackage{amsmath}
\title{Graph Search Algorithms}
\author{Kanishk Goel}
\date{14 November 2025}
\parindent 0px
\begin{document}
	\maketitle
	\section{Free Primitives}
	Treat any algorithm with linear running time as something we can essential use for free.
	\section{Generic Graph Search}
	\textbf{Input:}
	An undirected/directed graph $G = (V,E),$ and a starting vertex $s \in V$
	
	\textbf{Output:} Identify the vertices of $V$ reachable from $s$ in $G$. \\
	\begin{center}
		\begin{large}
			\textbf{Algorithm:} \\
		\end{large}
	\end{center}
	\begin{tabbing}
		\qquad \= \qquad \= \kill
		\textbf{Input:} graph $G = (V,E)$ and vertex $s \in V$ \\
		\textbf{Postcondition:} a vertex is reachable from $s$ if and only if it is marked as "explored" \\\\
		mark $s$ as explored and all others unexplored \\
		\textbf{while} $\exists\,(v,w)\in E$ with $v$ explored and $w$ unexplored \textbf{do} \\
		\> choose some edge $(v,w)$ \\
		\> mark $w$ as explored \\
		\textbf{end while}
	\end{tabbing}
	\section{Breadth First Search} 
	Explores layer by layer, Layer $0$ being $s$ itself.
	\begin{tabbing}
		\qquad \= \qquad \= \qquad \= \qquad \= \kill
		mark $s$ as explored, all other vertices as unexplored \\
		$Q$ := a queue data structure, init with $s$ \\
		while $Q$ is not empty do \\
		\> remove the vertex from the front of $Q$, call it $v$ \\
		\> for each edge $(v,w)$ in $v$'s adjacency list do \\
		\> \> if $w$ is unexplored then \\
		\> \> \> mark $w$ as explored \\
		\> \> \> add $w$ to the end of queue
	\end{tabbing} Runtime - $O(m+n)$
	
	\section{Connected Components}
	In an undirected graph $G = (V,E)$, a maximal subset $S \subseteq V$  of vertices such that there is a path from any vertex in $S$ to any other vertex in $S$.
	UCC algorithm is used which takes undirected graph $G = (V,E)$ in adjacency-list representation with $V = {1, 2, 3, \ldots, n}$ \\
	Postcondition: $\forall \: u,v \in V, \, cc(u) = cc(v) \iff u,v$ are in the same connected component
	\begin{center}
		\begin{large}
			\textbf{UCC}
		\end{large}
	\end{center}
		\begin{tabbing}
		\qquad \= \qquad \= \qquad \= \qquad \= \qquad \= \kill
		mark all vertices as unexplored \\
		numCC := 0 \\
		for i:= 1 to n do \\
		\> if i is unexplored then \\
		\> \> numCC := numCC + 1 //new component \\
		\> \> // call BFS \\
		\> \> $Q$ := a queue data structure, init with $s$ \\
		\> \> while $Q$ is not empty do \\
		\> \> \> remove the vertex from the front of $Q$, call it $v$ \\
		\> \> \> for each edge $(v,w)$ in $v$'s adjacency list do \\
		\> \> \> \> if $w$ is unexplored then \\
		\> \> \> \> \> mark $w$ as explored \\
		\> \> \>  \> \> add $w$ to the end of queue
	\end{tabbing} Runtime - $O(m+n)$
	
	\section{Depth First Search}
	Takes one path and only backtracks when absolutely necessary.
	
	\begin{center}
		\begin{large}
			\textbf{Iterative Implementation}
		\end{large}
	\end{center}
	
	\begin{tabbing}
		\qquad \= \qquad \= \qquad \= \qquad \= \kill
		mark all vertices as unexplored \\
		$S$ := a stack data structure, init with $s$ \\
		while $S$ is not empty do \\
		\> pop the vertex $v$ from the front of $S$ \\
		\> if $v$ is unexplored then \\
		\> \> mark $v$ as explored \\
		\> \> for each edge $(v,w)$ in $v$'s adjacency list do \\
		\> \> \> push $w$ to the front of $S$
	\end{tabbing}
	\pagebreak
	
	\begin{center}
		\begin{large}
			\textbf{Recursive Implementation}
		\end{large}
	\end{center}
	
	\begin{tabbing}
		\qquad \= \qquad \= \qquad \= \qquad \= \kill
		Func DFS($G,v$) \\
		\> mark $s$ as explored and all other as unexplored \\
		\> for each edge $(s,v)$ in $s$'s adjacency list do \\
		\> \> if $v$ is unexplored then\\
		\> \>  \> DFS($G,v$)
	\end{tabbing} Runtime - $O(m + n)$
	\section{Topological Sort}
	
	\begin{center}
		\begin{large}
			\textbf{Topological Ordering}
		\end{large}
	\end{center}
	\vspace{0.1mm}
	Let $G = (V,E)$ be a directed graph. A topological ordering of $G$ is an assignment $f(v)$ of every vertex $v \in V$ to a different number such that \[\forall \, (v,w) \in E, \: f(v) < f(w)\]
	
	\begin{itemize}
		\item Topological orderings do not exist in directed cyclic graphs.
		\item Every DAG has a topological ordering.
	\end{itemize}	
	
	\vspace{0.3mm}
	
	\begin{center}
		\begin{large}
			\textbf{Topological Sort}
		\end{large}
	\end{center}
	\vspace{0.1mm}
	DAG $G = (V,E)$ in adjacency list representation is given a $f$ value for every vertex constituting a topological ordering. it used DFS-Topo which assigns every vertex reachable from a vertex $s$ an $f$-value.
	\begin{tabbing}
		\qquad \= \qquad \= \qquad \= \kill
		Func TopoSort$(G)$ \\
		\> mark all vertices as unexplored \\
		\> curLabel := $|V|$ \\
		\> for every $v \in V$ do \\
		\> \> if $v$ is unexplored then \\
		\> \> \> DFS-Topo$(G,v)$
	\end{tabbing}
	
	\begin{tabbing}
		\qquad \= \qquad \= \qquad \= \kill
		Func DFS-Topo$(G,s)$ \\
		\> mark $s$ as explored \\
		\> for each edge $(s,v)$ in $s$'s outgoing adjacency list do \\
		\> \> if $v$ is unexplored then \\
		\> \> \> DFS-Topo$(G,v)$ \\
		\> $f(s)$ := curLabel \\
		\> curLabel := curLabel - 1
	\end{tabbing} curLabel is a global variable, initialised to $|V|$. Runtime is $O(m+n)$.
	
	\section{Strongly Connected Components}
	A strongly connected component or SCC of a directed graph is a maximal subset $S \subseteq V$ of vertices such that there is a directed path from any vertex in $S$ to any other vertex in $S$.
	
	Every directed graph can be viewed as directed acyclic graph built up from its SCCs.
	
	\subsection*{The SCC Meta-Graph is Directed Acyclic}
	$Let G = (V,E)$ be a directed graph. Define the corresponding meta-graph $H = (X,F)$ with one meta-vertex  $x \in X$ per SCC of $G$ and a meta-edge $(x,y)$ in $F$ whenever there is an edge in G from a vertex in the SCC corresponding to $x$ to one in the SCC corresponding to $y$. Then $H$ is a directed acyclic graph. 
	
	
\end{document}